% Осень26!A23; нумерация внутри столбца без пустых строк.
% Источник: лекция 3, §1 (строка 38)
% Источник: лекция 3, §2 (строка 58)
% Источник: лекция 3, §3 (строка 70)
\ticket{20}{Теорема о непрерывности меры сверху}
\label{ticket:20}

\begin{theorem}
Пусть $\mu$ --- мера на $\sigma$-алгебре $\mathcal A$, $A_n\in\mathcal A$, $A_1\supset A_2\supset\cdots$, $\mu(A_1)<\infty$.
Тогда
\[
  \mu\left(\bigcap_{n=1}^{\infty}A_n\right)
  =\lim_{n\to\infty}\mu(A_n).
\]
Обратно, объём $\mu$ на $\mathcal A$ с $\mu(X)<\infty$, обладающий этим свойством, является мерой.
\end{theorem}

\begin{proof}
Пусть $A=\bigcap_n A_n$ и $B_k=A_k\setminus A_{k+1}$.
По счётной аддитивности
\[
  \mu(A_1)=\mu(A)+\sum_{k=1}^{\infty}\mu(B_k)<\infty,
\]
\[
  \mu(A_n)=\mu(A)+\sum_{k=n}^{\infty}\mu(B_k).
\]
Остаток сходящегося ряда стремится к нулю, откуда следует первое утверждение.

Для обратного утверждения возьмём $C=\bigsqcup_{i=1}^{\infty}C_i$ и положим
\[
  R_n=C\setminus\bigcup_{i=1}^n C_i.
\]
Тогда $R_n\downarrow\varnothing$, поэтому $\mu(R_n)\to0$ по непрерывности сверху.
Конечная аддитивность даёт
\[
  \mu(C)=\sum_{i=1}^n\mu(C_i)+\mu(R_n).
\]
При $n\to\infty$ получаем счётную аддитивность.
\end{proof}
